% !TeX root = ../papers/nc-p3.tex
\section{半分の長さの三次曲線による三成分核の全面排除}
\label{sec:ternary-kernel-cubic-exclusion}
\begin{theorem}\label{thm:ternary-kernel-cubic-exclusion}
任意の四辺有限な標準有限Frobenius窓で、射影同値でない反対座標の対称二次式行列が
一般階数3なら、その核の原始ベクトルの四成分は線型独立である。
特に全局所環が曲線状の有限窓では、非自明な反対行列の行列式は必ず非零である。
局所長が全て偶数のときも例外はない。
\end{theorem}
\begin{proof}[平方零イデアルと三次曲線]
前節で残った三成分の核の場合を考える。左上三行 $H$ は一般階数2であり、
$\operatorname{adj}H=cgg^{\mathsf t}$ で $g$ は三つの原始な二次式である。
$G=(g_1,g_2,g_3)$ の $X$ 上の像を $J$ とする。
$J^2=0$ だから完全跡について $J$ は全等方である。
一方 $Y=X\cap V(G)$ の長さは前節により10以下で、$X$ の長さは20である。
従って $J$ も $Y$ も長さ10、各局所因子で $J=\operatorname{ann}J$、
$Y$ の局所長は $X$ の局所長のちょうど半分である。
特に全局所長は偶数である。この議論には曲線状性を使わない。

$H=P Q^{\mathsf t}+Q P^{\mathsf t}$ と書ける。実際、非零な主二次小行列式は定数倍の平方なので
分数体上で二次形式が二因子に分かれる。六つの独立な二次係数は原始で、
Gaussの補題により多項式因子を取れる。固定一次因子を持つ三変数二次形式の空間は三次元に過ぎず、
両因子の係数列 $P,Q$ は $x$ の一次式である。
$P\times Q$ は定数倍を除き $g$ で、$T=V(G)$ は前節の等号と
Hilbert--Burch分解による純次数3の曲線である。
また $N$ の第四列の先頭三成分 $b$ は $g$ のシジジーなので、
二つの最小線型シジジー $P,Q$ により $b=\ell P+mQ$（$\ell,m$ は一次式）と書ける。
最後の対角成分を $f$ とする。

$T$ の成分サイクルの全てが直線なら $Y$ の長さは高々9である。
ここで非被約な三次曲線にも使える長さ評価を記す。
純曲線の構造層で各既約被約成分 $E_i$ が滑らかな場合、一般重複度を $m_i$ とすると、
任意の有限部分スキーム $Y$ に対して
\[
              \operatorname{length}Y\le
              \sum_i m_i\operatorname{length}(Y\cap E_i).
\]
実際一つの一般点のArtin加群中で剰余体一本の部分加群を取り、
全一般点への局所化でこの部分加群に入る切断（他の一般点では零）を取る。
得た部分層は連接で、商には零次元部分層がなく、元の層の純性から部分層は
被約成分のイデアルで消える。その滑らかな成分上の捩れなし階数1層なので線束である。
これを繰り返せば各成分上の線束がちょうど $m_i$ 回現れるフィルトレーションを得る。
$\mathcal O_Y$ とのテンソルの右完全性で上の和が上界となる。
今、各直線と $X$ の切断長は3以下なので直線だけの場合を除外できる。

従って $T$ は非平面整三次、または非特異二次曲線と直線である。
前者が平面に入らないのは次数と純性により被約で、最小生成次数2から一次の方程式がないためである。
その正規化に次数3の線束の独立な四切断があるので、有理正規三次曲線となる。
後者では上の長さ評価により二次曲線と $Y$ の切断長は少なくとも7である。
この非線型成分を $E$、次数を $d=3,2$ とすると、どちらでも
$\operatorname{length}(Y\cap E)>2d+2$ となる。
\end{proof}

\begin{proof}[動く核点での次数と第四列]
三行四列の鉛筆 $uP+vQ$ の一般階数は3である。
一般階数が2ならその核直線が動いて曲面を掃くか、一つの固定直線が全核となる。
前者は $T$ の純一次元性に反し、後者では $P,Q$ が二つの変数だけを使うので $H$ の六成分が従属する。
階数1以下の鉛筆の元も、その核平面が $T$ に入るので存在しない。
従って非線型成分 $E$ は一般の鉛筆のただ一つの核点から得られる。
一定の鉛筆パラメータに属する成分は高々直線だからである。
比例係数から逆向きの有理写像があり、$E\simeq\mathbb P^1$ 上でこれは次数1のパラメータを与える。
$P,Q$ の共通零点はない。あれば一般の核点がそこに固定される。
従って $E$ 上では $P=a(s)V(s), Q=b_0(s)V(s)$ と書け、
$a,b_0$ は独立な二項一次式、$V$ の三成分は共通零点のない $d-1$ 次式である。
実際比例比の分母・分子を互いに素に取り、$P,Q$ の次数 $d$ の切断から割ればよい。

$N|_E$ の第四行列を含む主方向の二次小行列式は
\[
 V_iV_j\big(2ab_0 f-(a\ell+b_0m)^2\big).
\]
括弧は $2d+2$ 次で、$Y\cap E$ 上で零である。従って恒等零である。
$P,Q$ のうち少なくとも一方は階数3である。両方が階数2ならその核直線二本は
$T$ に含まれ、同じなら共通核があり、違えば非線型成分を入れると次数3を超える。
$P$ が階数3とし、その唯一の零点を $o$ とする。これは $E$ 上の $a=0$ の点であり
$b_0\ne0$。上の恒等式から $m(o)=0$、すなわち $m$ は $P$ の三成分の線型結合である。
第四行列への合同変換で $b$ から $Hc$ を引き $mQ$ の項を消すと、
$b=L P$（$L$ は一次式）となる。新しい最後の対角成分を再び $f$ と書く。
すると $X$ 上で
\[
 P_i u_j+P_j u_i=0,\qquad
          u_i=fQ_i-\tfrac12 L^2 P_i .
\]
$o$ が $X$ の台にあれば $f(o)\ne0$ である。$N(o)$ の他の成分は全て零だが
全二次式が同時に消える点はないからである。$Q(o)\ne0$ ゆえその局所環で一つの $u_i$ が単元となる。
対角式から $P_i=0$、残りの式から全 $P_j=0$ となり、この点は長さ1である。
全局所長が偶数であることに反する。よって $P$ は $X$ 上で基点を持たず、各点で単元な $P_i$ を使うと
逆に全 $u_j=0$ となる。三次評価の単射性により $fQ=\tfrac12L^2P$ が多項式として成立する。
これは $P\times Q\ne0$ と $f\ne0$ に反する。
曲線状の場合の四成分核は前節の重み評価と整閉性ですでに全て除かれている。
\end{proof}
